Um zu zeigen, dass $\operatorname{UT}_n(\mathbb{K}) = \operatorname{SUT}_n(\mathbb{K}) \rtimes \operatorname{Diag}_n(\mathbb{K})$, betrachten wir die Struktur der beteiligten Gruppen und die Bedingungen für das semidirekte Produkt.

Die Gruppe $\operatorname{UT}_n(\mathbb{K})$ besteht aus allen oberen Dreiecksmatrizen mit Einträgen aus einem Körper $\mathbb{K}$ und nicht verschwindenden Diagonaleinträgen. Die Gruppe $\operatorname{SUT}_n(\mathbb{K})$ ist die Untergruppe von $\operatorname{UT}_n(\mathbb{K})$, die aus oberen Dreiecksmatrizen mit Einsen auf der Diagonale besteht. Schließlich ist $\operatorname{Diag}_n(\mathbb{K})$ die Gruppe der diagonalen Matrizen mit nicht verschwindenden Einträgen aus $\mathbb{K}$.

Um die semidirekte Produktstruktur zu zeigen, müssen wir zwei Dinge nachweisen: Erstens, dass $\operatorname{SUT}_n(\mathbb{K})$ ein Normalteiler von $\operatorname{UT}_n(\mathbb{K})$ ist, und zweitens, dass jede Matrix in $\operatorname{UT}_n(\mathbb{K})$ eindeutig als Produkt einer Matrix aus $\operatorname{SUT}_n(\mathbb{K})$ und einer Matrix aus $\operatorname{Diag}_n(\mathbb{K})$ geschrieben werden kann.

Zunächst zeigen wir die Normalität von $\operatorname{SUT}_n(\mathbb{K})$ in $\operatorname{UT}_n(\mathbb{K})$. Sei $A \in \operatorname{SUT}_n(\mathbb{K})$ und $D \in \operatorname{Diag}_n(\mathbb{K})$. Wir betrachten das Produkt $DAD^{-1}$:

\begin{align*}
D &= \operatorname{diag}(d_1, d_2, \ldots, d_n), \\
D^{-1} &= \operatorname{diag}(d_1^{-1}, d_2^{-1}, \ldots, d_n^{-1}).
\end{align*}

Dann ist

\begin{align*}
DAD^{-1} &= \begin{pmatrix}
d_1 & 0 & \cdots & 0 \\
0 & d_2 & \cdots & 0 \\
\vdots & \vdots & \ddots & \vdots \\
0 & 0 & \cdots & d_n
\end{pmatrix}
\begin{pmatrix}
1 & a_{12} & \cdots & a_{1n} \\
0 & 1 & \cdots & a_{2n} \\
\vdots & \vdots & \ddots & \vdots \\
0 & 0 & \cdots & 1
\end{pmatrix}
\begin{pmatrix}
d_1^{-1} & 0 & \cdots & 0 \\
0 & d_2^{-1} & \cdots & 0 \\
\vdots & \vdots & \ddots & \vdots \\
0 & 0 & \cdots & d_n^{-1}
\end{pmatrix} \\
&= \begin{pmatrix}
1 & d_1 a_{12} d_2^{-1} & \cdots & d_1 a_{1n} d_n^{-1} \\
0 & 1 & \cdots & d_2 a_{2n} d_n^{-1} \\
\vdots & \vdots & \ddots & \vdots \\
0 & 0 & \cdots & 1
\end{pmatrix}.
\end{align*}

Da die Diagonalelemente von $DAD^{-1}$ Einsen sind, folgt $DAD^{-1} \in \operatorname{SUT}_n(\mathbb{K})$, was die Normalität zeigt.

Nun zur Eindeutigkeit der Darstellung. Sei $U \in \operatorname{UT}_n(\mathbb{K})$. Dann kann $U$ als Produkt $U = SD$ geschrieben werden, wobei $S \in \operatorname{SUT}_n(\mathbb{K})$ und $D \in \operatorname{Diag}_n(\mathbb{K})$. Die Matrix $D$ wird gebildet, indem man die Diagonalelemente von $U$ nimmt, und $S$ erhält man, indem man $U$ durch $D$ von rechts multipliziert, also $S = UD^{-1}$. Da $D$ eindeutig durch die Diagonalelemente von $U$ bestimmt ist, ist auch $S$ eindeutig bestimmt.

Daher ist $\operatorname{UT}_n(\mathbb{K})$ tatsächlich das semidirekte Produkt von $\operatorname{SUT}_n(\mathbb{K})$ und $\operatorname{Diag}_n(\mathbb{K})$, was die Behauptung beweist.

%CHECK: Explizite Berechnung von DAD^{-1} zur Bestätigung der Normalität hinzugefügt.